\section{Read a modular inverse from the same divisions}
An inverse of \(e\) modulo \(m\) is an integer \(d\) for which
\(ed\equiv1\pmod m\). Such an inverse exists precisely when \(e\) and \(m\) are
coprime. The gcd calculation tells us whether that condition holds. Its divisions can
also recover coefficients \(u\) and \(v\) satisfying \(um+ve=1\).

Work backwards from the remainder one in the preceding example:
\[
\begin{aligned}
1&=9-8,\\
 &=9-(17-9)=2\cdot9-17,\\
 &=2(3120-183\cdot17)-17=2\cdot3120-367\cdot17.
\end{aligned}
\]
Consequently \(-367\) is an inverse of 17 modulo 3120. Add 3120 to obtain the
nonnegative representative 2753. Verify it directly:
\(17\cdot2753=46801=15\cdot3120+1\).

For a toy RSA example, let \(p=61\), \(q=53\), \(n=pq=3233\) and
\(\varphi(n)=(p-1)(q-1)=3120\). The public exponent 17 and private exponent 2753
satisfy this inverse relation. We will use those fixed values to exercise modular
power. This is not a request to write a key generator: no primality testing or random
prime selection is supplied by the arithmetic interfaces.

Keep the inverse calculation on paper for this project. Extending Euclid to calculate
coefficients in assembly introduces signed values and coefficient bounds, which need
a separate representation contract. The unsigned gcd return alone does not encode
those coefficients. In particular, the negative coefficient \(-367\) is not an unsigned
128-bit input to our power function.
